\section{Teknologi yang Digunakan}

Selama pelaksanaan magang, berbagai teknologi digunakan untuk mendukung proses
pengembangan perangkat lunak, kolaborasi tim, serta aktivitas operasional
sehari-hari. Teknologi-teknologi tersebut dapat dikelompokkan sebagai berikut:

\begin{enumerate}
      \item \textbf{Framework dan bahasa pemrograman}

            \begin{enumerate}[noitemsep, topsep=0pt]
                  \item \textbf{Flutter (Dart)}

                        Flutter merupakan \textit{framework} UI yang dikembangkan oleh Google untuk
                        membangun aplikasi \textit{multiplatform} menggunakan satu \textit{codebase}
                        \cite{flutter_docs}. Pada periode awal magang (Januari--Februari 2026), Flutter
                        digunakan dalam pengembangan dan pemeliharaan aplikasi \textit{mobile}
                        HealthyOne dan Canvasser. Implementasi dilakukan menggunakan pendekatan
                        \textit{Server-Driven UI} (SDUI), yaitu pendekatan di mana struktur antarmuka
                        sebagian ditentukan dari sisi server melalui konfigurasi yang dikirimkan ke
                        aplikasi \cite{airbnb_sd_ui}.

                  \item \textbf{React (Next.js dan TypeScript)}

                        React merupakan \textit{library} JavaScript yang dikembangkan oleh Meta untuk
                        membangun UI berbasis komponen \cite{react_docs}. Dalam implementasinya, React
                        dipadukan dengan Next.js yang menyediakan fitur \textit{routing}, optimasi, dan
                        mekanisme \textit{server-side rendering} \cite{nextjs_docs}, serta TypeScript
                        yang memberikan dukungan sistem tipe statis untuk meningkatkan kualitas dan
                        keterbacaan kode. Kombinasi teknologi ini digunakan mulai pertengahan Februari
                        2026 sebagai \textit{tech stack} utama pada pengembangan ulang UI HealthyOne
                        dan Canvasser.

                  \item \textbf{Toraja (Django/Python)}

                        Toraja merupakan \textit{framework} internal perusahaan yang dibangun
                        menggunakan Django dan Python. Django sendiri merupakan \textit{framework} web
                        tingkat tinggi yang menerapkan prinsip pengembangan cepat dan desain yang
                        terstruktur \cite{django_docs}. Pada periode pertengahan April 2026, Toraja
                        digunakan untuk membuat dan menyesuaikan \textit{endpoint} API yang
                        menghubungkan antarmuka React dengan database seiring dengan berkembangnya
                        peran menjadi \textit{Full-Stack Developer}.
            \end{enumerate}

      \item \textbf{\textit{Library} dan \textit{tools} pendukung}

            Pada pengembangan aplikasi berbasis React, digunakan beberapa \textit{library}
            dan \textit{tool} pendukung. Mantine UI menyediakan komponen UI yang
            \textit{reusable}, sedangkan Tailwind CSS digunakan untuk mempermudah proses
            \textit{styling}. Selain itu, AG Grid digunakan sebagai dasar implementasi
            komponen DataView yang mendukung pengelolaan data berbentuk tabel secara
            interaktif \cite{aggrid_docs}. Pada sisi arsitektur aplikasi, digunakan lapisan
            \textit{services}, \textit{hooks}, \textit{mapper}, \textit{helper}, dan
            \textit{utils} untuk mendukung pengelolaan \textit{state} serta komunikasi API
            dengan Toraja.

      \item \textbf{Perangkat kolaborasi dan produktivitas}

            Selain teknologi pengembangan aplikasi, digunakan pula beberapa perangkat
            pendukung untuk menunjang kolaborasi tim. OpenProject dimanfaatkan sebagai
            sarana manajemen \textit{task} dan pencatatan KPI, sedangkan GitLab digunakan
            sebagai media \textit{version control} serta proses \textit{merge request}.
            Untuk mendukung aktivitas operasional perusahaan, digunakan Sistem Informasi
            Pelaporan Perusahaan (SIPP) Pharos untuk presensi dan administrasi karyawan,
            serta WhatsApp sebagai media komunikasi antar anggota tim.
\end{enumerate}
